\chapter{Analisi dei requisiti}
\label{chap:requisiti}

% ==================================================
% 3.1 Analisi del problema
% ==================================================

\section{Analisi del problema}
\label{sec:analisi-problema}

% Obiettivo:
% tradurre il problema introdotto nei Capitoli 1 e 2 in
% esigenze progettuali verificabili.
%
% Contenuti attesi:
% - attori principali del sistema;
% - patrimonio librario privato come risorsa gestita;
% - discovery bibliografica e territoriale;
% - vincoli di privacy sulla posizione;
% - interazione tra proprietario e utente interessato.
%
% Non descrivere ancora l'architettura tecnica: le scelte
% implementative appartengono ai Capitoli 4 e 5.

% ==================================================
% 3.2 Attori e casi d'uso
% ==================================================

\section{Attori e casi d'uso}
\label{sec:attori-casi-uso}

% Attori previsti:
% - visitatore non autenticato;
% - utente registrato;
% - proprietario di libri;
% - amministratore.
%
% Casi d'uso da derivare dal prototipo reale:
% - registrazione e conferma email;
% - autenticazione;
% - gestione profilo;
% - pubblicazione, modifica ed eliminazione libro;
% - consultazione catalogo pubblico;
% - ricerca geografica protetta;
% - invio e gestione di richieste;
% - consultazione dashboard e statistiche;
% - accesso al pannello amministrativo.

% ==================================================
% 3.3 Requisiti funzionali
% ==================================================

\section{Requisiti funzionali}
\label{sec:requisiti-funzionali}

% I requisiti devono essere numerati in modo stabile
% (RF-01, RF-02, ...), collegabili ai casi d'uso e ai test
% del Capitolo 6.
%
% Scrivere solo requisiti supportati dal codice corrente o
% chiaramente identificati come pianificati.
%
% Aree funzionali:
% - account e autenticazione;
% - catalogazione libraria;
% - immagini e copertine;
% - catalogazione assistita da ISBN/OCR;
% - ricerca e discovery;
% - ricerca per prossimità;
% - richieste di consultazione/prestito;
% - dashboard utente;
% - amministrazione.

% ==================================================
% 3.4 Requisiti non funzionali
% ==================================================

\section{Requisiti non funzionali}
\label{sec:requisiti-non-funzionali}

% Requisiti non funzionali da trattare:
% - privacy della posizione;
% - sicurezza dell'autenticazione;
% - modularità e manutenibilità;
% - riproducibilità dell'ambiente;
% - accessibilità e responsività essenziali;
% - internazionalizzazione;
% - affidabilità dei fallback verso provider esterni.
%
% Evitare requisiti non misurabili o non verificabili.

% ==================================================
% 3.5 Vincoli e assunzioni
% ==================================================

\section{Vincoli e assunzioni progettuali}
\label{sec:vincoli-assunzioni}

% Descrivere i vincoli che influenzano la progettazione:
% - prototipo universitario;
% - esecuzione locale e servizi configurabili;
% - uso di provider esterni opzionali;
% - necessità di non esporre coordinate precise;
% - limiti della validazione manuale rispetto a una suite
%   automatizzata completa.
